\scanpage{79}
\mspage{II. 27}
% archive sheet number: 78
% line 1
4. \uline{Un complément au th. principal}.
% lines 2–6
\uline{Proposition} 4.1 Sous les conditions de 1.1, pour tout \deleted{\unsure{$\ell$}} nb premier $\ell$ sauf un nb fini, \deleted{le groupe} \deleted{\unsure{$\pi$}} on a $\overline{[\pi',I(\ell)]}=I(\ell)$, i.e. le plus grand quotient \unsure{profini|propre} de $I(\ell)$ sur lequel $\pi_0$ opère trivialement est \uline{nul}.
% lines 7–10
On voit par dualité, en procédant comme dans 2.4, que cet énoncé équivaut au suivant :

\deleted{\uline{Proposition}} \inserted{Corollaire} 4.2. Pour presque tout $\ell$, on a
\[H^1(\bar X,\mathbb Z/\ell\mathbb Z)^{\pi_0}=0.\]
% lines 11–17
Pour prouver ceci, on est encore ramené, \deleted{\unsure{comme}} \deleted{\unsure{au cas où}} \deleted{\unsure{qu'on a}} par la réduction de 2.3, au cas où \deleted{$X\subset X'$} $X$ est un ouvert d'un $X'$ projectif et géom normal sur $k$. Dans ce cas, \deleted{le lemme} 3.2. nous ramène au cas où de plus $X'$ est de dim 1, donc est une courbe lisse et complète.
% lines 18–21
Quitte à \deleted{\unsure{agrandir $k$}} passer à une extension finie de $k$, on peut encore supposer les pts \inserted{$x_i$} de $X'-X$ rationnels sur $k$. On utilise alors la suite exacte analogue à (2.5.1)
% lines 22–24
\[(4.2.1)\]
\diagram{\begin{tikzcd}[ampersand replacement=\&,column sep=small]
{} \arrow[r,dash] \& H^1(\bar X',\mathbb Z/\ell\mathbb Z) \arrow[r] \& H^1(\bar X,\mathbb Z/\ell\mathbb Z) \arrow[r] \& \deleted{\unsure{H}}\prod_i H^2_x(\bar X,\mathbb Z/\ell\mathbb Z) \arrow[r,dash] \& H^2(\bar X',\unsure{\mathbb Z(\ell)|\mathbb Z/\ell\mathbb Z}) \arrow[d,dash]\\
\&\&\&\&{}
\end{tikzcd}}
% lines 25–28
qui nous montre que $H^1(\bar X,\mathbb Z/\ell\mathbb Z)$ est une extension de $\bigl(\mu_\ell(\bar k^*)^{\otimes-1}\bigr)^{n-1}$ par $H^1(\bar X',\mathbb Z/\ell\mathbb Z)$.

\deleted{\unsure{Comme le gr}} Soit
\[N=\text{card du grp des racines de 1 dans }k^*.\tag*{(4.2.2)}\]
